\paragraph{What the strategy comparison shows.}
\begin{itemize}
  \item \textbf{Warm start is the largest single gain.} Without it, DQN with the new reward (R2)
        waits 308\,s at real demand and 945\,s at 1.3$\times$ demand, close to the fixed plan.
        With the camera-rule warm start (R3--R5) the mean is 109\,s and 622\,s.
  \item \textbf{A longer view helps the learned controller.} R10 (warm start, full 236\,m view)
        waits 76\,s against 86\,s for actuated control (11\% less) and 467\,s against 491\,s at
        1.3$\times$ demand (5\% less). One camera at the stop line cannot give this view; it needs
        a second camera up the road or a higher mounting, so R10 is an upper bound and not part of
        the camera-only comparison.
  \item \textbf{The hand-written camera rule is the strongest controller in this test}
        (66\,s and 453\,s). It already beats actuated control, because it skips empty roads and
        ranks roads by PCU queue and waiting time. DQN starts from this rule but does not keep all
        of its quality; with a camera-only view the three seeds differ widely (96\,s, 87\,s and
        146\,s at real demand).
  \item \textbf{Longer training is worse, not better.} R8 (300{,}000 steps) waits 431\,s against
        109\,s for the same setup at 150{,}000 steps. The agent drifts away from the teacher's
        policy as training continues, which points to a need for early stopping or a stronger
        anchor to the teacher.
  \item \textbf{The other changes do not help in this setting}: the keep/switch action (R6),
        QR-DQN (R7) and the 5\,s step (R1). R1 has a low waiting time at 1.3$\times$ demand only
        because 37\% of its vehicles never entered the road and so have no waiting time recorded;
        its throughput is the lowest of all controllers.
  \item \textbf{Reward:} the combined reward gives a lower waiting time than the old reward at real
        demand (109\,s against 150\,s, R3--R5 against R9) but a higher one at 1.3$\times$ demand
        (622\,s against 549\,s), mostly because of one weak seed.
\end{itemize}
The full test therefore uses the camera rule and the camera-view DQN (the R3--R5 setup), which
both work with one camera at the stop line. Its result in the full test is the mean of several training seeds.
